% Generated by reports/ycb1_tools/ycb1_make_book_figures.py from docs/image_register.csv. Do not edit by hand.
% One prompt and one final file per figure, used by \ycbfig{<figure_id>} (see template_book.tex).

\expandafter\def\csname figprompt@ycb1-fig-autonomic-nervous-system\endcsname{Anatomy diagram: Autonomic Nervous System (ANS). A simple schematic in front or side view, whichever shows the structure best; main parts as flat shapes with thin outlines, grayscale-safe. Labels: only the terms listed under this figure in the notes block (list them here before generating); each label short, horizontal, with a leader line. One idea only; leave out finer detail. Aspect ratio 4:3.}
\expandafter\def\csname figfile@ycb1-fig-autonomic-nervous-system\endcsname{../workshop/images/new/Autonomic_Nervous_System_new.pdf}
\expandafter\def\csname figprompt@ycb1-fig-basic-brain-anatomy\endcsname{Anatomy diagram: Central Nervous System: Brain. A simple schematic in front or side view, whichever shows the structure best; main parts as flat shapes with thin outlines, grayscale-safe. Labels: only the terms listed under this figure in the notes block (list them here before generating); each label short, horizontal, with a leader line. One idea only; leave out finer detail. Aspect ratio 4:3.}
\expandafter\def\csname figfile@ycb1-fig-basic-brain-anatomy\endcsname{../workshop/images/new/Basic-brain-anatomy_new.pdf}
\expandafter\def\csname figprompt@ycb1-fig-cardiovascularsystem\endcsname{Anatomy diagram: Cardiovascular System. A simple schematic in front or side view, whichever shows the structure best; main parts as flat shapes with thin outlines, grayscale-safe. Labels: only the terms listed under this figure in the notes block (list them here before generating); each label short, horizontal, with a leader line. One idea only; leave out finer detail. Aspect ratio 4:3.}
\expandafter\def\csname figfile@ycb1-fig-cardiovascularsystem\endcsname{../workshop/images/new/CardiovascularSystem_new.pdf}
\expandafter\def\csname figprompt@ycb1-fig-heart-diagram\endcsname{Anatomy diagram: Blood Vessels and Heart. A simple schematic in front or side view, whichever shows the structure best; main parts as flat shapes with thin outlines, grayscale-safe. Labels: only the terms listed under this figure in the notes block (list them here before generating); each label short, horizontal, with a leader line. One idea only; leave out finer detail. Aspect ratio 4:3.}
\expandafter\def\csname figfile@ycb1-fig-heart-diagram\endcsname{../workshop/images/new/Heart-diagram_new.pdf}
\expandafter\def\csname figprompt@ycb1-fig-human-excretory-system\endcsname{Anatomy diagram: Excretory System. A simple schematic in front or side view, whichever shows the structure best; main parts as flat shapes with thin outlines, grayscale-safe. Labels: only the terms listed under this figure in the notes block (list them here before generating); each label short, horizontal, with a leader line. One idea only; leave out finer detail. Aspect ratio 4:3.}
\expandafter\def\csname figfile@ycb1-fig-human-excretory-system\endcsname{../workshop/images/new/Human-Excretory-System_new.pdf}
\expandafter\def\csname figprompt@ycb1-fig-human-male-reproductive-system-en\endcsname{Anatomy diagram: Male Reproductive System. A simple schematic in front or side view, whichever shows the structure best; main parts as flat shapes with thin outlines, grayscale-safe. Labels: only the terms listed under this figure in the notes block (list them here before generating); each label short, horizontal, with a leader line. One idea only; leave out finer detail. Aspect ratio 4:3.}
\expandafter\def\csname figfile@ycb1-fig-human-male-reproductive-system-en\endcsname{../workshop/images/new/Human_male_reproductive_system_en_new.pdf}
\expandafter\def\csname figprompt@ycb1-fig-mechnaism-of-muscle-contrtaction\endcsname{Anatomy diagram: Muscle Contraction. A simple schematic in front or side view, whichever shows the structure best; main parts as flat shapes with thin outlines, grayscale-safe. Labels: only the terms listed under this figure in the notes block (list them here before generating); each label short, horizontal, with a leader line. One idea only; leave out finer detail. Aspect ratio 4:3.}
\expandafter\def\csname figfile@ycb1-fig-mechnaism-of-muscle-contrtaction\endcsname{../workshop/images/new/Mechnaism-of-muscle-contrtaction_new.pdf}
\expandafter\def\csname figprompt@ycb1-fig-muscle-chart-showing-the-muscular-system-labeled\endcsname{Anatomy diagram: Muscular System Overview. A simple schematic in front or side view, whichever shows the structure best; main parts as flat shapes with thin outlines, grayscale-safe. Labels: only the terms listed under this figure in the notes block (list them here before generating); each label short, horizontal, with a leader line. One idea only; leave out finer detail. Aspect ratio 4:3.}
\expandafter\def\csname figfile@ycb1-fig-muscle-chart-showing-the-muscular-system-labeled\endcsname{../workshop/images/new/Muscle-chart-showing-the-muscular-system-labeled_new.pdf}
\expandafter\def\csname figprompt@ycb1-fig-peripheral-nervous-system\endcsname{Anatomy diagram: Peripheral Nervous System: Overview. A simple schematic in front or side view, whichever shows the structure best; main parts as flat shapes with thin outlines, grayscale-safe. Labels: only the terms listed under this figure in the notes block (list them here before generating); each label short, horizontal, with a leader line. One idea only; leave out finer detail. Aspect ratio 4:3.}
\expandafter\def\csname figfile@ycb1-fig-peripheral-nervous-system\endcsname{../workshop/images/new/Peripheral_Nervous_System_new.pdf}
\expandafter\def\csname figprompt@ycb1-fig-reciprocal-inhibition\endcsname{Anatomy diagram: Reflex Action and Reciprocal Inhibition. A simple schematic in front or side view, whichever shows the structure best; main parts as flat shapes with thin outlines, grayscale-safe. Labels: only the terms listed under this figure in the notes block (list them here before generating); each label short, horizontal, with a leader line. One idea only; leave out finer detail. Aspect ratio 4:3.}
\expandafter\def\csname figfile@ycb1-fig-reciprocal-inhibition\endcsname{../workshop/images/new/Reciprocal-Inhibition_new.pdf}
\expandafter\def\csname figprompt@ycb1-fig-somatic-nervous-system\endcsname{Anatomy diagram: Somatic Nervous System (SNS). A simple schematic in front or side view, whichever shows the structure best; main parts as flat shapes with thin outlines, grayscale-safe. Labels: only the terms listed under this figure in the notes block (list them here before generating); each label short, horizontal, with a leader line. One idea only; leave out finer detail. Aspect ratio 4:3.}
\expandafter\def\csname figfile@ycb1-fig-somatic-nervous-system\endcsname{../workshop/images/new/Somatic_nervous_system_new.pdf}
\expandafter\def\csname figprompt@ycb1-fig-spinal-chord\endcsname{Anatomy diagram: Central Nervous System: Spinal Cord. A simple schematic in front or side view, whichever shows the structure best; main parts as flat shapes with thin outlines, grayscale-safe. Labels: only the terms listed under this figure in the notes block (list them here before generating); each label short, horizontal, with a leader line. One idea only; leave out finer detail. Aspect ratio 4:3.}
\expandafter\def\csname figfile@ycb1-fig-spinal-chord\endcsname{../workshop/images/new/Spinal_chord_new.pdf}
\expandafter\def\csname figprompt@ycb1-fig-article-respiration-muscles-wc1\endcsname{Anatomy diagram: Muscle Breathing. A simple schematic in front or side view, whichever shows the structure best; main parts as flat shapes with thin outlines, grayscale-safe. Labels: only the terms listed under this figure in the notes block (list them here before generating); each label short, horizontal, with a leader line. One idea only; leave out finer detail. Aspect ratio 4:3.}
\expandafter\def\csname figfile@ycb1-fig-article-respiration-muscles-wc1\endcsname{../workshop/images/new/article-Respiration-Muscles-WC1_new.pdf}
\expandafter\def\csname figprompt@ycb1-fig-breathing\endcsname{Anatomy diagram: Muscles of Respiration. A simple schematic in front or side view, whichever shows the structure best; main parts as flat shapes with thin outlines, grayscale-safe. Labels: only the terms listed under this figure in the notes block (list them here before generating); each label short, horizontal, with a leader line. One idea only; leave out finer detail. Aspect ratio 4:3.}
\expandafter\def\csname figfile@ycb1-fig-breathing\endcsname{../workshop/images/new/breathing_new.pdf}
\expandafter\def\csname figprompt@ycb1-fig-digestive-system\endcsname{Anatomy diagram: Digestive System. A simple schematic in front or side view, whichever shows the structure best; main parts as flat shapes with thin outlines, grayscale-safe. Labels: only the terms listed under this figure in the notes block (list them here before generating); each label short, horizontal, with a leader line. One idea only; leave out finer detail. Aspect ratio 4:3.}
\expandafter\def\csname figfile@ycb1-fig-digestive-system\endcsname{../workshop/images/new/digestive-system_new.pdf}
\expandafter\def\csname figprompt@ycb1-fig-endocrine-system\endcsname{Anatomy diagram: Endocrine System. A simple schematic in front or side view, whichever shows the structure best; main parts as flat shapes with thin outlines, grayscale-safe. Labels: only the terms listed under this figure in the notes block (list them here before generating); each label short, horizontal, with a leader line. One idea only; leave out finer detail. Aspect ratio 4:3.}
\expandafter\def\csname figfile@ycb1-fig-endocrine-system\endcsname{../workshop/images/new/endocrine-system_new.pdf}
\expandafter\def\csname figprompt@ycb1-fig-endocrine\endcsname{Anatomy diagram: Endocrine Glands: Pineal, Thyroid, and Parathyroid. A simple schematic in front or side view, whichever shows the structure best; main parts as flat shapes with thin outlines, grayscale-safe. Labels: only the terms listed under this figure in the notes block (list them here before generating); each label short, horizontal, with a leader line. One idea only; leave out finer detail. Aspect ratio 4:3.}
\expandafter\def\csname figfile@ycb1-fig-endocrine\endcsname{../workshop/images/new/endocrine_new.pdf}
\expandafter\def\csname figprompt@ycb1-fig-female-reproductive-system\endcsname{Anatomy diagram: Female Reproductive System. A simple schematic in front or side view, whichever shows the structure best; main parts as flat shapes with thin outlines, grayscale-safe. Labels: only the terms listed under this figure in the notes block (list them here before generating); each label short, horizontal, with a leader line. One idea only; leave out finer detail. Aspect ratio 4:3.}
\expandafter\def\csname figfile@ycb1-fig-female-reproductive-system\endcsname{../workshop/images/new/female-reproductive-system_new.pdf}
\expandafter\def\csname figprompt@ycb1-fig-joints\endcsname{Anatomy diagram: Types of Joints. A simple schematic in front or side view, whichever shows the structure best; main parts as flat shapes with thin outlines, grayscale-safe. Labels: only the terms listed under this figure in the notes block (list them here before generating); each label short, horizontal, with a leader line. One idea only; leave out finer detail. Aspect ratio 4:3.}
\expandafter\def\csname figfile@ycb1-fig-joints\endcsname{../workshop/images/new/joints_new.pdf}
\expandafter\def\csname figprompt@ycb1-fig-overview-nervious-system\endcsname{Anatomy diagram: Overview of Nervous System. A simple schematic in front or side view, whichever shows the structure best; main parts as flat shapes with thin outlines, grayscale-safe. Labels: only the terms listed under this figure in the notes block (list them here before generating); each label short, horizontal, with a leader line. One idea only; leave out finer detail. Aspect ratio 4:3.}
\expandafter\def\csname figfile@ycb1-fig-overview-nervious-system\endcsname{../workshop/images/new/overview_nervious_system_new.pdf}
\expandafter\def\csname figprompt@ycb1-fig-pituitary-gland\endcsname{Anatomy diagram: Endocrine Glands: Hypothalamus and Pituitary. A simple schematic in front or side view, whichever shows the structure best; main parts as flat shapes with thin outlines, grayscale-safe. Labels: only the terms listed under this figure in the notes block (list them here before generating); each label short, horizontal, with a leader line. One idea only; leave out finer detail. Aspect ratio 4:3.}
\expandafter\def\csname figfile@ycb1-fig-pituitary-gland\endcsname{../workshop/images/new/pituitary-gland_new.pdf}
\expandafter\def\csname figprompt@ycb1-fig-reprosys\endcsname{Anatomy diagram: Overview of Reproductive System. A simple schematic in front or side view, whichever shows the structure best; main parts as flat shapes with thin outlines, grayscale-safe. Labels: only the terms listed under this figure in the notes block (list them here before generating); each label short, horizontal, with a leader line. One idea only; leave out finer detail. Aspect ratio 4:3.}
\expandafter\def\csname figfile@ycb1-fig-reprosys\endcsname{../workshop/images/new/reprosys_new.pdf}
\expandafter\def\csname figprompt@ycb1-fig-respiratory-system\endcsname{Anatomy diagram: Respiratory System. A simple schematic in front or side view, whichever shows the structure best; main parts as flat shapes with thin outlines, grayscale-safe. Labels: only the terms listed under this figure in the notes block (list them here before generating); each label short, horizontal, with a leader line. One idea only; leave out finer detail. Aspect ratio 4:3.}
\expandafter\def\csname figfile@ycb1-fig-respiratory-system\endcsname{../workshop/images/new/respiratory-system_new.pdf}
\expandafter\def\csname figprompt@ycb1-fig-skeleton\endcsname{Anatomy diagram: Components of the Skeletal System. A simple schematic in front or side view, whichever shows the structure best; main parts as flat shapes with thin outlines, grayscale-safe. Labels: only the terms listed under this figure in the notes block (list them here before generating); each label short, horizontal, with a leader line. One idea only; leave out finer detail. Aspect ratio 4:3.}
\expandafter\def\csname figfile@ycb1-fig-skeleton\endcsname{../workshop/images/new/skeleton_new.pdf}
\expandafter\def\csname figprompt@ycb1-fig-skeleton-divisions\endcsname{Anatomy diagram: Divisions of the Skeletal System. A simple schematic in front or side view, whichever shows the structure best; main parts as flat shapes with thin outlines, grayscale-safe. Labels: only the terms listed under this figure in the notes block (list them here before generating); each label short, horizontal, with a leader line. One idea only; leave out finer detail. Aspect ratio 4:3.}
\expandafter\def\csname figfile@ycb1-fig-skeleton-divisions\endcsname{../workshop/images/new/skeleton_divisions_new.pdf}
\expandafter\def\csname figprompt@ycb1-fig-spinal-movements\endcsname{Anatomy diagram: Types of Spinal Movements. A simple schematic in front or side view, whichever shows the structure best; main parts as flat shapes with thin outlines, grayscale-safe. Labels: only the terms listed under this figure in the notes block (list them here before generating); each label short, horizontal, with a leader line. One idea only; leave out finer detail. Aspect ratio 4:3.}
\expandafter\def\csname figfile@ycb1-fig-spinal-movements\endcsname{../workshop/images/new/spinal_movements_new.pdf}
\expandafter\def\csname figprompt@ycb1-fig-vertebral-column\endcsname{Anatomy diagram: Vertebral Column. A simple schematic in front or side view, whichever shows the structure best; main parts as flat shapes with thin outlines, grayscale-safe. Labels: only the terms listed under this figure in the notes block (list them here before generating); each label short, horizontal, with a leader line. One idea only; leave out finer detail. Aspect ratio 4:3.}
\expandafter\def\csname figfile@ycb1-fig-vertebral-column\endcsname{../workshop/images/new/vertebral_column_new.pdf}
\expandafter\def\csname figprompt@ycb1-fig-abdominalbreathing\endcsname{Concept diagram: Abdominal Breathing (Diaphragmatic Breathing). Boxes, circles and arrows only, drawn from the structure given in the notes block (the names and their order are listed there); Sanskrit terms in Roman letters with the English meaning in brackets, nothing else added. Flat colours that still read in grayscale. Best redrawn later in TikZ. Aspect ratio 4:3.}
\expandafter\def\csname figfile@ycb1-fig-abdominalbreathing\endcsname{../workshop/images/new/ycb1_abdominalbreathing_new.pdf}
\expandafter\def\csname figprompt@ycb1-fig-abductors\endcsname{Muscle action: Abductors. Simple stick figure in side view, two overlapping positions (start and end) in light and dark line, the working muscle group drawn as a thick coloured line along the limb, one thin arrow for the movement. Label only with the single word 'Abductors'. Aspect ratio 4:3.}
\expandafter\def\csname figfile@ycb1-fig-abductors\endcsname{../workshop/images/new/ycb1_abductors_new.pdf}
\expandafter\def\csname figprompt@ycb1-fig-adductors\endcsname{Muscle action: Adductors. Simple stick figure in side view, two overlapping positions (start and end) in light and dark line, the working muscle group drawn as a thick coloured line along the limb, one thin arrow for the movement. Label only with the single word 'Adductors'. Aspect ratio 4:3.}
\expandafter\def\csname figfile@ycb1-fig-adductors\endcsname{../workshop/images/new/ycb1_adductors_new.pdf}
\expandafter\def\csname figprompt@ycb1-fig-anulomvilom\endcsname{Pranayama: Anulom Vilom (Alternate Nostril Breathing). One adult practitioner (CS1), seated in a comfortable cross-legged pose, spine upright, eyes closed. Hands, fingers, face and tongue or lips exactly as in the notes block for this practice; mark the side of the breath with thin arrows only if the notes block gives a side. Front three-quarter view, head and shoulders large enough to show the hand and face. Aspect ratio 4:5 portrait.}
\expandafter\def\csname figfile@ycb1-fig-anulomvilom\endcsname{../workshop/images/new/ycb1_anulomvilom_new.png}
\expandafter\def\csname figprompt@ycb1-fig-application-of-knowledge\endcsname{Concept diagram: Application of Knowledge as Yoga Protocol Instructor. Boxes, circles and arrows only, drawn from the structure given in the notes block (the names and their order are listed there); Sanskrit terms in Roman letters with the English meaning in brackets, nothing else added. Flat colours that still read in grayscale. Best redrawn later in TikZ. Aspect ratio 4:3.}
\expandafter\def\csname figfile@ycb1-fig-application-of-knowledge\endcsname{../workshop/images/new/ycb1_application_of_knowledge_new.pdf}
\expandafter\def\csname figprompt@ycb1-fig-ardhachakrasan\endcsname{Posture: Ardha Chakrasana or Hastottanasan. One adult practitioner (CS1) in the final held position of Ardha Chakrasana or Hastottanasan, side view, standing, facing right. The whole body in frame with a small margin; head, spine, arms, hands, legs and feet positions exactly as in the notes block for this posture; weight and support points clear. No props unless the notes block names one. No labels. Aspect ratio 4:5 portrait.}
\expandafter\def\csname figfile@ycb1-fig-ardhachakrasan\endcsname{../workshop/images/new/ycb1_ardhachakrasan_new.png}
\expandafter\def\csname figprompt@ycb1-fig-asanabenefits\endcsname{Do not generate a picture. This is a screenshot of a table. Retype it as a LaTeX table (deck) and a Markdown table (notes) from the checked notes block ycb1-pd-yogasana-practice, in own words where it is a list of facts.}
\expandafter\def\csname figfile@ycb1-fig-asanabenefits\endcsname{../workshop/images/new/}
\expandafter\def\csname figprompt@ycb1-fig-asanabenefits-1\endcsname{Do not generate a picture. This is a screenshot of a table. Retype it as a LaTeX table (deck) and a Markdown table (notes) from the checked notes block ycb1-pd-yogasana-practice, in own words where it is a list of facts.}
\expandafter\def\csname figfile@ycb1-fig-asanabenefits-1\endcsname{../workshop/images/new/}
\expandafter\def\csname figprompt@ycb1-fig-asanabenefits-2\endcsname{Do not generate a picture. This is a screenshot of a table. Retype it as a LaTeX table (deck) and a Markdown table (notes) from the checked notes block ycb1-pd-yogasana-practice, in own words where it is a list of facts.}
\expandafter\def\csname figfile@ycb1-fig-asanabenefits-2\endcsname{../workshop/images/new/}
\expandafter\def\csname figprompt@ycb1-fig-asanabenefits-3\endcsname{Do not generate a picture. This is a screenshot of a table. Retype it as a LaTeX table (deck) and a Markdown table (notes) from the checked notes block ycb1-pd-yogasana-practice, in own words where it is a list of facts.}
\expandafter\def\csname figfile@ycb1-fig-asanabenefits-3\endcsname{../workshop/images/new/}
\expandafter\def\csname figprompt@ycb1-fig-bandhas\endcsname{Concept diagram: Bandhas (Internal Locks). Boxes, circles and arrows only, drawn from the structure given in the notes block (the names and their order are listed there); Sanskrit terms in Roman letters with the English meaning in brackets, nothing else added. Flat colours that still read in grayscale. Best redrawn later in TikZ. Aspect ratio 4:3.}
\expandafter\def\csname figfile@ycb1-fig-bandhas\endcsname{../workshop/images/new/ycb1_bandhas_new.pdf}
\expandafter\def\csname figprompt@ycb1-fig-basti\endcsname{Cleansing technique: Basti. A clean schematic only: the practitioner (CS1) in the starting posture of the technique with a small simple inset showing the path or the area involved, as a flat outline with no internal organs drawn in detail and nothing graphic. Show the posture and the direction of the action with thin arrows. No tools or fluids drawn realistically. Aspect ratio 4:3. This technique is taught only under a qualified teacher; the picture is a memory aid.}
\expandafter\def\csname figfile@ycb1-fig-basti\endcsname{../workshop/images/new/ycb1_basti_new.png}
\expandafter\def\csname figprompt@ycb1-fig-bhadrasan\endcsname{Posture: Bhadrasana. One adult practitioner (CS1) in the final held position of Bhadrasana, front three-quarter view, seated on a plain mat. The whole body in frame with a small margin; head, spine, arms, hands, legs and feet positions exactly as in the notes block for this posture; weight and support points clear. No props unless the notes block names one. No labels. Aspect ratio 4:5 portrait.}
\expandafter\def\csname figfile@ycb1-fig-bhadrasan\endcsname{../workshop/images/new/ycb1_bhadrasan_new.png}
\expandafter\def\csname figprompt@ycb1-fig-bhramari\endcsname{Pranayama: Bhramari (Bee Breath). One adult practitioner (CS1), seated in a comfortable cross-legged pose, spine upright, eyes closed. Hands, fingers, face and tongue or lips exactly as in the notes block for this practice; mark the side of the breath with thin arrows only if the notes block gives a side. Front three-quarter view, head and shoulders large enough to show the hand and face. Aspect ratio 4:5 portrait.}
\expandafter\def\csname figfile@ycb1-fig-bhramari\endcsname{../workshop/images/new/ycb1_bhramari_new.png}
\expandafter\def\csname figprompt@ycb1-fig-bhuja\endcsname{Movement practice: Shoulder Movement. One adult practitioner (CS1), standing, front or side view as the joint needs. Show the start position and the end position of the movement in one picture, as two overlapping figures in light and normal tone, with one thin curved arrow for the direction. The moving joint is clearly visible; the rest of the body stays still. No props. Aspect ratio 3:2 landscape.}
\expandafter\def\csname figfile@ycb1-fig-bhuja\endcsname{../workshop/images/new/ycb1_bhuja_new.png}
\expandafter\def\csname figprompt@ycb1-fig-bhujangasan\endcsname{Posture: Bhujangasana. One adult practitioner (CS1) in the final held position of Bhujangasana, side view, lying face down on a plain mat. The whole body in frame with a small margin; head, spine, arms, hands, legs and feet positions exactly as in the notes block for this posture; weight and support points clear. No props unless the notes block names one. No labels. Aspect ratio 3:2 landscape.}
\expandafter\def\csname figfile@ycb1-fig-bhujangasan\endcsname{../workshop/images/new/ycb1_bhujangasan_new.png}
\expandafter\def\csname figprompt@ycb1-fig-bodysystems\endcsname{Do not generate a picture. This is a screenshot of a table. Retype it as a LaTeX table (deck) and a Markdown table (notes) from the checked notes block ycb1-hl-human-body-introduction, in own words where it is a list of facts.}
\expandafter\def\csname figfile@ycb1-fig-bodysystems\endcsname{../workshop/images/new/}
\expandafter\def\csname figprompt@ycb1-fig-chakras\endcsname{Concept diagram: Chakra Mudras (Energy Center Gestures); Chakras. Boxes, circles and arrows only, drawn from the structure given in the notes block (the names and their order are listed there); Sanskrit terms in Roman letters with the English meaning in brackets, nothing else added. Flat colours that still read in grayscale. Best redrawn later in TikZ. Aspect ratio 4:3.}
\expandafter\def\csname figfile@ycb1-fig-chakras\endcsname{../workshop/images/new/ycb1_chakras_new.pdf}
\expandafter\def\csname figprompt@ycb1-fig-chestbreathing\endcsname{Concept diagram: Chest Breathing. Boxes, circles and arrows only, drawn from the structure given in the notes block (the names and their order are listed there); Sanskrit terms in Roman letters with the English meaning in brackets, nothing else added. Flat colours that still read in grayscale. Best redrawn later in TikZ. Aspect ratio 4:3.}
\expandafter\def\csname figfile@ycb1-fig-chestbreathing\endcsname{../workshop/images/new/ycb1_chestbreathing_new.pdf}
\expandafter\def\csname figprompt@ycb1-fig-clavicularbreathing\endcsname{Preparatory breathing: Clavicular Breathing. One adult practitioner (CS1), side view, seated or lying as the notes block says, hands placed as the notes block says. Show the torso with a soft translucent highlight on the part that rises on the in-breath and a thin arrow for the direction of movement; a faint outline of the lungs and diaphragm may sit behind the highlight. Calm face, eyes closed. Aspect ratio 4:3 landscape.}
\expandafter\def\csname figfile@ycb1-fig-clavicularbreathing\endcsname{../workshop/images/new/ycb1_clavicularbreathing_new.png}
\expandafter\def\csname figprompt@ycb1-fig-dandasan\endcsname{Posture: Dandasana. One adult practitioner (CS1) in the final held position of Dandasana, front three-quarter view, seated on a plain mat. The whole body in frame with a small margin; head, spine, arms, hands, legs and feet positions exactly as in the notes block for this posture; weight and support points clear. No props unless the notes block names one. No labels. Aspect ratio 4:5 portrait.}
\expandafter\def\csname figfile@ycb1-fig-dandasan\endcsname{../workshop/images/new/ycb1_dandasan_new.png}
\expandafter\def\csname figprompt@ycb1-fig-dhatus\endcsname{Concept diagram: Dhatus. Boxes, circles and arrows only, drawn from the structure given in the notes block (the names and their order are listed there); Sanskrit terms in Roman letters with the English meaning in brackets, nothing else added. Flat colours that still read in grayscale. Best redrawn later in TikZ. Aspect ratio 4:3.}
\expandafter\def\csname figfile@ycb1-fig-dhatus\endcsname{../workshop/images/new/ycb1_dhatus_new.pdf}
\expandafter\def\csname figprompt@ycb1-fig-dhauti\endcsname{Cleansing technique: Dhauti. A clean schematic only: the practitioner (CS1) in the starting posture of the technique with a small simple inset showing the path or the area involved, as a flat outline with no internal organs drawn in detail and nothing graphic. Show the posture and the direction of the action with thin arrows. No tools or fluids drawn realistically. Aspect ratio 4:3. This technique is taught only under a qualified teacher; the picture is a memory aid.}
\expandafter\def\csname figfile@ycb1-fig-dhauti\endcsname{../workshop/images/new/ycb1_dhauti_new.png}
\expandafter\def\csname figprompt@ycb1-fig-exam-i\endcsname{Do not generate a picture. This is a screenshot of a table. Retype it as a LaTeX table (deck) and a Markdown table (notes) from the checked notes block ycb1-ex-exam-and-certification, in own words where it is a list of facts.}
\expandafter\def\csname figfile@ycb1-fig-exam-i\endcsname{../workshop/images/new/}
\expandafter\def\csname figprompt@ycb1-fig-exam-ii\endcsname{Do not generate a picture. This is a screenshot of a table. Retype it as a LaTeX table (deck) and a Markdown table (notes) from the checked notes block ycb1-ex-exam-and-certification, in own words where it is a list of facts.}
\expandafter\def\csname figfile@ycb1-fig-exam-ii\endcsname{../workshop/images/new/}
\expandafter\def\csname figprompt@ycb1-fig-exam-iii\endcsname{Do not generate a picture. This is a screenshot of a table. Retype it as a LaTeX table (deck) and a Markdown table (notes) from the checked notes block ycb1-ex-exam-and-certification, in own words where it is a list of facts.}
\expandafter\def\csname figfile@ycb1-fig-exam-iii\endcsname{../workshop/images/new/}
\expandafter\def\csname figprompt@ycb1-fig-extensors\endcsname{Muscle action: Extensors. Simple stick figure in side view, two overlapping positions (start and end) in light and dark line, the working muscle group drawn as a thick coloured line along the limb, one thin arrow for the movement. Label only with the single word 'Extensors'. Aspect ratio 4:3.}
\expandafter\def\csname figfile@ycb1-fig-extensors\endcsname{../workshop/images/new/ycb1_extensors_new.pdf}
\expandafter\def\csname figprompt@ycb1-fig-flexors\endcsname{Muscle action: Flexors. Simple stick figure in side view, two overlapping positions (start and end) in light and dark line, the working muscle group drawn as a thick coloured line along the limb, one thin arrow for the movement. Label only with the single word 'Flexors'. Aspect ratio 4:3.}
\expandafter\def\csname figfile@ycb1-fig-flexors\endcsname{../workshop/images/new/ycb1_flexors_new.pdf}
\expandafter\def\csname figprompt@ycb1-fig-gomukhasan\endcsname{Posture: Gomukhasana. One adult practitioner (CS1) in the final held position of Gomukhasana, front three-quarter view, seated on a plain mat. The whole body in frame with a small margin; head, spine, arms, hands, legs and feet positions exactly as in the notes block for this posture; weight and support points clear. No props unless the notes block names one. No labels. Aspect ratio 4:5 portrait.}
\expandafter\def\csname figfile@ycb1-fig-gomukhasan\endcsname{../workshop/images/new/ycb1_gomukhasan_new.png}
\expandafter\def\csname figprompt@ycb1-fig-greeva\endcsname{Movement practice: Neck Movement. One adult practitioner (CS1), standing, front or side view as the joint needs. Show the start position and the end position of the movement in one picture, as two overlapping figures in light and normal tone, with one thin curved arrow for the direction. The moving joint is clearly visible; the rest of the body stays still. No props. Aspect ratio 3:2 landscape.}
\expandafter\def\csname figfile@ycb1-fig-greeva\endcsname{../workshop/images/new/ycb1_greeva_new.png}
\expandafter\def\csname figprompt@ycb1-fig-greevabend\endcsname{Movement practice: Neck Movement. One adult practitioner (CS1), standing, front or side view as the joint needs. Show the start position and the end position of the movement in one picture, as two overlapping figures in light and normal tone, with one thin curved arrow for the direction. The moving joint is clearly visible; the rest of the body stays still. No props. Aspect ratio 3:2 landscape.}
\expandafter\def\csname figfile@ycb1-fig-greevabend\endcsname{../workshop/images/new/ycb1_greevabend_new.png}
\expandafter\def\csname figprompt@ycb1-fig-gulpha\endcsname{Movement practice: Ankle Movement. One adult practitioner (CS1), standing, front or side view as the joint needs. Show the start position and the end position of the movement in one picture, as two overlapping figures in light and normal tone, with one thin curved arrow for the direction. The moving joint is clearly visible; the rest of the body stays still. No props. Aspect ratio 3:2 landscape.}
\expandafter\def\csname figfile@ycb1-fig-gulpha\endcsname{../workshop/images/new/ycb1_gulpha_new.png}
\expandafter\def\csname figprompt@ycb1-fig-hastamudra\endcsname{Mudra practice: Hasta Mudras (Hand Gestures). Close-up of the hand or hands in the gesture named in the notes block, one hand per panel if more than one gesture is shown, same skin tone in every panel, five fingers per hand, fingertips and thumb contact clearly drawn. Small inset of the seated practitioner (CS1) showing where the hands rest. Aspect ratio 4:3.}
\expandafter\def\csname figfile@ycb1-fig-hastamudra\endcsname{../workshop/images/new/ycb1_hastamudra_new.png}
\expandafter\def\csname figprompt@ycb1-fig-history\endcsname{Concept diagram: History of Yoga. Boxes, circles and arrows only, drawn from the structure given in the notes block (the names and their order are listed there); Sanskrit terms in Roman letters with the English meaning in brackets, nothing else added. Flat colours that still read in grayscale. Best redrawn later in TikZ. Aspect ratio 4:3.}
\expandafter\def\csname figfile@ycb1-fig-history\endcsname{../workshop/images/new/ycb1_history_new.pdf}
\expandafter\def\csname figprompt@ycb1-fig-hymns\endcsname{Concept diagram: Concept of Hymns. Boxes, circles and arrows only, drawn from the structure given in the notes block (the names and their order are listed there); Sanskrit terms in Roman letters with the English meaning in brackets, nothing else added. Flat colours that still read in grayscale. Best redrawn later in TikZ. Aspect ratio 4:3.}
\expandafter\def\csname figfile@ycb1-fig-hymns\endcsname{../workshop/images/new/ycb1_hymns_new.pdf}
\expandafter\def\csname figprompt@ycb1-fig-jangha\endcsname{Movement practice: Jangha. One adult practitioner (CS1), standing, front or side view as the joint needs. Show the start position and the end position of the movement in one picture, as two overlapping figures in light and normal tone, with one thin curved arrow for the direction. The moving joint is clearly visible; the rest of the body stays still. No props. Aspect ratio 3:2 landscape.}
\expandafter\def\csname figfile@ycb1-fig-jangha\endcsname{../workshop/images/new/ycb1_jangha_new.png}
\expandafter\def\csname figprompt@ycb1-fig-jangha-1\endcsname{Movement practice: Thighs Movement. One adult practitioner (CS1), standing, front or side view as the joint needs. Show the start position and the end position of the movement in one picture, as two overlapping figures in light and normal tone, with one thin curved arrow for the direction. The moving joint is clearly visible; the rest of the body stays still. No props. Aspect ratio 3:2 landscape.}
\expandafter\def\csname figfile@ycb1-fig-jangha-1\endcsname{../workshop/images/new/ycb1_jangha_1_new.png}
\expandafter\def\csname figprompt@ycb1-fig-jangha-2\endcsname{Movement practice: Thighs Movement. One adult practitioner (CS1), standing, front or side view as the joint needs. Show the start position and the end position of the movement in one picture, as two overlapping figures in light and normal tone, with one thin curved arrow for the direction. The moving joint is clearly visible; the rest of the body stays still. No props. Aspect ratio 3:2 landscape.}
\expandafter\def\csname figfile@ycb1-fig-jangha-2\endcsname{../workshop/images/new/ycb1_jangha_2_new.png}
\expandafter\def\csname figprompt@ycb1-fig-janu-1\endcsname{Movement practice: Knee Movement. One adult practitioner (CS1), standing, front or side view as the joint needs. Show the start position and the end position of the movement in one picture, as two overlapping figures in light and normal tone, with one thin curved arrow for the direction. The moving joint is clearly visible; the rest of the body stays still. No props. Aspect ratio 3:2 landscape.}
\expandafter\def\csname figfile@ycb1-fig-janu-1\endcsname{../workshop/images/new/ycb1_janu_1_new.png}
\expandafter\def\csname figprompt@ycb1-fig-janu-2\endcsname{Concept illustration: Knee Movement. One practitioner (CS1) performing the practice described in the notes block, plus one small clean schematic if the concept needs it, nothing else. Aspect ratio 4:3.}
\expandafter\def\csname figfile@ycb1-fig-janu-2\endcsname{../workshop/images/new/ycb1_janu_2_new.png}
\expandafter\def\csname figprompt@ycb1-fig-kapalbhati\endcsname{Practice: Kapalabhati (Skull Shining Breath); Kapalabhati. One adult practitioner (CS1), seated, front three-quarter view, in the posture of the notes block. Show the key action in a small inset: the abdomen drawn in sharply for kapalabhati, the throat narrowed softly for ujjayi, the steady gaze on a candle flame at eye level for trataka, the tongue position for the mukh mudra. Use only the inset that matches Kapalabhati (Skull Shining Breath); Kapalabhati. Aspect ratio 4:5 portrait.}
\expandafter\def\csname figfile@ycb1-fig-kapalbhati\endcsname{../workshop/images/new/ycb1_kapalbhati_new.png}
\expandafter\def\csname figprompt@ycb1-fig-kati\endcsname{Movement practice: Trunk Movement. One adult practitioner (CS1), standing, front or side view as the joint needs. Show the start position and the end position of the movement in one picture, as two overlapping figures in light and normal tone, with one thin curved arrow for the direction. The moving joint is clearly visible; the rest of the body stays still. No props. Aspect ratio 3:2 landscape.}
\expandafter\def\csname figfile@ycb1-fig-kati\endcsname{../workshop/images/new/ycb1_kati_new.png}
\expandafter\def\csname figprompt@ycb1-fig-katichakrasan\endcsname{Posture: Kati Chakrasana. One adult practitioner (CS1) in the final held position of Kati Chakrasana, side view, standing, facing right. The whole body in frame with a small margin; head, spine, arms, hands, legs and feet positions exactly as in the notes block for this posture; weight and support points clear. No props unless the notes block names one. No labels. Aspect ratio 4:5 portrait.}
\expandafter\def\csname figfile@ycb1-fig-katichakrasan\endcsname{../workshop/images/new/ycb1_katichakrasan_new.png}
\expandafter\def\csname figprompt@ycb1-fig-kayamudra\endcsname{Mudra practice: Kaya Mudras (Body Gestures). Close-up of the hand or hands in the gesture named in the notes block, one hand per panel if more than one gesture is shown, same skin tone in every panel, five fingers per hand, fingertips and thumb contact clearly drawn. Small inset of the seated practitioner (CS1) showing where the hands rest. Aspect ratio 4:3.}
\expandafter\def\csname figfile@ycb1-fig-kayamudra\endcsname{../workshop/images/new/ycb1_kayamudra_new.png}
\expandafter\def\csname figprompt@ycb1-fig-khechari\endcsname{Practice: Mukh Mudras (Facial Gestures). One adult practitioner (CS1), seated, front three-quarter view, in the posture of the notes block. Show the key action in a small inset: the abdomen drawn in sharply for kapalabhati, the throat narrowed softly for ujjayi, the steady gaze on a candle flame at eye level for trataka, the tongue position for the mukh mudra. Use only the inset that matches Mukh Mudras (Facial Gestures). Aspect ratio 4:5 portrait.}
\expandafter\def\csname figfile@ycb1-fig-khechari\endcsname{../workshop/images/new/ycb1_khechari_new.png}
\expandafter\def\csname figprompt@ycb1-fig-kosha\endcsname{Concept diagram: (Five Sheaths). Boxes, circles and arrows only, drawn from the structure given in the notes block (the names and their order are listed there); Sanskrit terms in Roman letters with the English meaning in brackets, nothing else added. Flat colours that still read in grayscale. Best redrawn later in TikZ. Aspect ratio 4:3.}
\expandafter\def\csname figfile@ycb1-fig-kosha\endcsname{../workshop/images/new/ycb1_kosha_new.pdf}
\expandafter\def\csname figprompt@ycb1-fig-makarasan\endcsname{Posture: Makarasana. One adult practitioner (CS1) in the final held position of Makarasana, side view, lying face down on a plain mat. The whole body in frame with a small margin; head, spine, arms, hands, legs and feet positions exactly as in the notes block for this posture; weight and support points clear. No props unless the notes block names one. No labels. Aspect ratio 3:2 landscape.}
\expandafter\def\csname figfile@ycb1-fig-makarasan\endcsname{../workshop/images/new/ycb1_makarasan_new.png}
\expandafter\def\csname figprompt@ycb1-fig-mandukasan\endcsname{Posture: Mandukasana. One adult practitioner (CS1) in the final held position of Mandukasana, front three-quarter view, seated on a plain mat. The whole body in frame with a small margin; head, spine, arms, hands, legs and feet positions exactly as in the notes block for this posture; weight and support points clear. No props unless the notes block names one. No labels. Aspect ratio 4:5 portrait.}
\expandafter\def\csname figfile@ycb1-fig-mandukasan\endcsname{../workshop/images/new/ycb1_mandukasan_new.png}
\expandafter\def\csname figprompt@ycb1-fig-nauli\endcsname{Cleansing technique: Nauli. A clean schematic only: the practitioner (CS1) in the starting posture of the technique with a small simple inset showing the path or the area involved, as a flat outline with no internal organs drawn in detail and nothing graphic. Show the posture and the direction of the action with thin arrows. No tools or fluids drawn realistically. Aspect ratio 4:3. This technique is taught only under a qualified teacher; the picture is a memory aid.}
\expandafter\def\csname figfile@ycb1-fig-nauli\endcsname{../workshop/images/new/ycb1_nauli_new.png}
\expandafter\def\csname figprompt@ycb1-fig-neti\endcsname{Cleansing technique: Neti. A clean schematic only: the practitioner (CS1) in the starting posture of the technique with a small simple inset showing the path or the area involved, as a flat outline with no internal organs drawn in detail and nothing graphic. Show the posture and the direction of the action with thin arrows. No tools or fluids drawn realistically. Aspect ratio 4:3. This technique is taught only under a qualified teacher; the picture is a memory aid.}
\expandafter\def\csname figfile@ycb1-fig-neti\endcsname{../workshop/images/new/ycb1_neti_new.png}
\expandafter\def\csname figprompt@ycb1-fig-pada\endcsname{Movement practice: Ankle Movement. One adult practitioner (CS1), standing, front or side view as the joint needs. Show the start position and the end position of the movement in one picture, as two overlapping figures in light and normal tone, with one thin curved arrow for the direction. The moving joint is clearly visible; the rest of the body stays still. No props. Aspect ratio 3:2 landscape.}
\expandafter\def\csname figfile@ycb1-fig-pada\endcsname{../workshop/images/new/ycb1_pada_new.png}
\expandafter\def\csname figprompt@ycb1-fig-padahastasan\endcsname{Posture: Padahastasana. One adult practitioner (CS1) in the final held position of Padahastasana, side view, standing, facing right. The whole body in frame with a small margin; head, spine, arms, hands, legs and feet positions exactly as in the notes block for this posture; weight and support points clear. No props unless the notes block names one. No labels. Aspect ratio 4:5 portrait.}
\expandafter\def\csname figfile@ycb1-fig-padahastasan\endcsname{../workshop/images/new/ycb1_padahastasan_new.png}
\expandafter\def\csname figprompt@ycb1-fig-padmasan\endcsname{Posture: Padmasana. One adult practitioner (CS1) in the final held position of Padmasana, front three-quarter view, seated on a plain mat. The whole body in frame with a small margin; head, spine, arms, hands, legs and feet positions exactly as in the notes block for this posture; weight and support points clear. No props unless the notes block names one. No labels. Aspect ratio 4:5 portrait.}
\expandafter\def\csname figfile@ycb1-fig-padmasan\endcsname{../workshop/images/new/ycb1_padmasan_new.png}
\expandafter\def\csname figprompt@ycb1-fig-pashchimottansasan\endcsname{Posture: Paschimottanasana. One adult practitioner (CS1) in the final held position of Paschimottanasana, front three-quarter view, seated on a plain mat. The whole body in frame with a small margin; head, spine, arms, hands, legs and feet positions exactly as in the notes block for this posture; weight and support points clear. No props unless the notes block names one. No labels. Aspect ratio 4:5 portrait.}
\expandafter\def\csname figfile@ycb1-fig-pashchimottansasan\endcsname{../workshop/images/new/ycb1_pashchimottansasan_new.png}
\expandafter\def\csname figprompt@ycb1-fig-pavanamuktasan\endcsname{Posture: Pavanamuktasana. One adult practitioner (CS1) in the final held position of Pavanamuktasana, side view, lying on the back on a plain mat. The whole body in frame with a small margin; head, spine, arms, hands, legs and feet positions exactly as in the notes block for this posture; weight and support points clear. No props unless the notes block names one. No labels. Aspect ratio 3:2 landscape.}
\expandafter\def\csname figfile@ycb1-fig-pavanamuktasan\endcsname{../workshop/images/new/ycb1_pavanamuktasan_new.png}
\expandafter\def\csname figprompt@ycb1-fig-pranas\endcsname{Concept diagram: Pranas. Boxes, circles and arrows only, drawn from the structure given in the notes block (the names and their order are listed there); Sanskrit terms in Roman letters with the English meaning in brackets, nothing else added. Flat colours that still read in grayscale. Best redrawn later in TikZ. Aspect ratio 4:3.}
\expandafter\def\csname figfile@ycb1-fig-pranas\endcsname{../workshop/images/new/ycb1_pranas_new.pdf}
\expandafter\def\csname figprompt@ycb1-fig-pranav\endcsname{Concept diagram: Concept of Pranava. Boxes, circles and arrows only, drawn from the structure given in the notes block (the names and their order are listed there); Sanskrit terms in Roman letters with the English meaning in brackets, nothing else added. Flat colours that still read in grayscale. Best redrawn later in TikZ. Aspect ratio 4:3.}
\expandafter\def\csname figfile@ycb1-fig-pranav\endcsname{../workshop/images/new/ycb1_pranav_new.pdf}
\expandafter\def\csname figprompt@ycb1-fig-purvottansasan\endcsname{Posture: Purvottanasana. One adult practitioner (CS1) in the final held position of Purvottanasana, front three-quarter view, seated on a plain mat. The whole body in frame with a small margin; head, spine, arms, hands, legs and feet positions exactly as in the notes block for this posture; weight and support points clear. No props unless the notes block names one. No labels. Aspect ratio 4:5 portrait.}
\expandafter\def\csname figfile@ycb1-fig-purvottansasan\endcsname{../workshop/images/new/ycb1_purvottansasan_new.png}
\expandafter\def\csname figprompt@ycb1-fig-rotators\endcsname{Muscle action: Rotators. Simple stick figure in side view, two overlapping positions (start and end) in light and dark line, the working muscle group drawn as a thick coloured line along the limb, one thin arrow for the movement. Label only with the single word 'Rotators'. Aspect ratio 4:3.}
\expandafter\def\csname figfile@ycb1-fig-rotators\endcsname{../workshop/images/new/ycb1_rotators_new.pdf}
\expandafter\def\csname figprompt@ycb1-fig-saralmatsyasan\endcsname{Posture: Saral Matsyasana. One adult practitioner (CS1) in the final held position of Saral Matsyasana, side view, lying on the back on a plain mat. The whole body in frame with a small margin; head, spine, arms, hands, legs and feet positions exactly as in the notes block for this posture; weight and support points clear. No props unless the notes block names one. No labels. Aspect ratio 3:2 landscape.}
\expandafter\def\csname figfile@ycb1-fig-saralmatsyasan\endcsname{../workshop/images/new/ycb1_saralmatsyasan_new.png}
\expandafter\def\csname figprompt@ycb1-fig-setubandhasan\endcsname{Posture: Setubandhasana. One adult practitioner (CS1) in the final held position of Setubandhasana, side view, lying on the back on a plain mat. The whole body in frame with a small margin; head, spine, arms, hands, legs and feet positions exactly as in the notes block for this posture; weight and support points clear. No props unless the notes block names one. No labels. Aspect ratio 3:2 landscape.}
\expandafter\def\csname figfile@ycb1-fig-setubandhasan\endcsname{../workshop/images/new/ycb1_setubandhasan_new.png}
\expandafter\def\csname figprompt@ycb1-fig-shalabhasan\endcsname{Posture: Shalabhasana. One adult practitioner (CS1) in the final held position of Shalabhasana, side view, lying face down on a plain mat. The whole body in frame with a small margin; head, spine, arms, hands, legs and feet positions exactly as in the notes block for this posture; weight and support points clear. No props unless the notes block names one. No labels. Aspect ratio 3:2 landscape.}
\expandafter\def\csname figfile@ycb1-fig-shalabhasan\endcsname{../workshop/images/new/ycb1_shalabhasan_new.png}
\expandafter\def\csname figprompt@ycb1-fig-shashankasan\endcsname{Posture: Shashankasana. One adult practitioner (CS1) in the final held position of Shashankasana, front three-quarter view, seated on a plain mat. The whole body in frame with a small margin; head, spine, arms, hands, legs and feet positions exactly as in the notes block for this posture; weight and support points clear. No props unless the notes block names one. No labels. Aspect ratio 4:5 portrait.}
\expandafter\def\csname figfile@ycb1-fig-shashankasan\endcsname{../workshop/images/new/ycb1_shashankasan_new.png}
\expandafter\def\csname figprompt@ycb1-fig-sitali\endcsname{Pranayama: Sitali (Cooling Breath). One adult practitioner (CS1), seated in a comfortable cross-legged pose, spine upright, eyes closed. Hands, fingers, face and tongue or lips exactly as in the notes block for this practice; mark the side of the breath with thin arrows only if the notes block gives a side. Front three-quarter view, head and shoulders large enough to show the hand and face. Aspect ratio 4:5 portrait.}
\expandafter\def\csname figfile@ycb1-fig-sitali\endcsname{../workshop/images/new/ycb1_sitali_new.png}
\expandafter\def\csname figprompt@ycb1-fig-sthula\endcsname{Concept illustration: Sthula Vyayama: Concept. One practitioner (CS1) performing the practice described in the notes block, plus one small clean schematic if the concept needs it, nothing else. Aspect ratio 4:3.}
\expandafter\def\csname figfile@ycb1-fig-sthula\endcsname{../workshop/images/new/ycb1_sthula_new.png}
\expandafter\def\csname figprompt@ycb1-fig-sthula-hrid\endcsname{Posture: Hrid Gati (Engine Daud). One adult practitioner (CS1) in the final held position of Hrid Gati (Engine Daud), side view. The whole body in frame with a small margin; head, spine, arms, hands, legs and feet positions exactly as in the notes block for this posture; weight and support points clear. No props unless the notes block names one. No labels. Aspect ratio 4:5 portrait.}
\expandafter\def\csname figfile@ycb1-fig-sthula-hrid\endcsname{../workshop/images/new/ycb1_sthula_hrid_new.png}
\expandafter\def\csname figprompt@ycb1-fig-sthula-sarvanga\endcsname{Posture: Sarvanga Pushti. One adult practitioner (CS1) in the final held position of Sarvanga Pushti, side view. The whole body in frame with a small margin; head, spine, arms, hands, legs and feet positions exactly as in the notes block for this posture; weight and support points clear. No props unless the notes block names one. No labels. Aspect ratio 4:5 portrait.}
\expandafter\def\csname figfile@ycb1-fig-sthula-sarvanga\endcsname{../workshop/images/new/ycb1_sthula_sarvanga_new.png}
\expandafter\def\csname figprompt@ycb1-fig-sukhasan\endcsname{Posture: Sukhasana. One adult practitioner (CS1) in the final held position of Sukhasana, front three-quarter view, seated on a plain mat. The whole body in frame with a small margin; head, spine, arms, hands, legs and feet positions exactly as in the notes block for this posture; weight and support points clear. No props unless the notes block names one. No labels. Aspect ratio 4:5 portrait.}
\expandafter\def\csname figfile@ycb1-fig-sukhasan\endcsname{../workshop/images/new/ycb1_sukhasan_new.png}
\expandafter\def\csname figprompt@ycb1-fig-sukshma\endcsname{Concept illustration: Sukshma Vyayama: Concept. One practitioner (CS1) performing the practice described in the notes block, plus one small clean schematic if the concept needs it, nothing else. Aspect ratio 4:3.}
\expandafter\def\csname figfile@ycb1-fig-sukshma\endcsname{../workshop/images/new/ycb1_sukshma_new.png}
\expandafter\def\csname figprompt@ycb1-fig-sunsal1\endcsname{Do not generate a picture. This is a screenshot of a table. Retype it as a LaTeX table (deck) and a Markdown table (notes) from the checked notes block ycb1-pd-surya-namaskara, in own words where it is a list of facts.}
\expandafter\def\csname figfile@ycb1-fig-sunsal1\endcsname{../workshop/images/new/}
\expandafter\def\csname figprompt@ycb1-fig-sunsal2\endcsname{Do not generate a picture. This is a screenshot of a table. Retype it as a LaTeX table (deck) and a Markdown table (notes) from the checked notes block ycb1-pd-surya-namaskara, in own words where it is a list of facts.}
\expandafter\def\csname figfile@ycb1-fig-sunsal2\endcsname{../workshop/images/new/}
\expandafter\def\csname figprompt@ycb1-fig-suryanamaskar\endcsname{The Surya Namaskara sequence: one practitioner (CS1) in a strip of panels, one panel per position in the order given in the notes block, side view, each panel with a thin arrow to the next and a breath marker (in or out) as the notes block says. Panels equal size, white gaps, no text except the breath marker symbols. Aspect ratio 16:9.}
\expandafter\def\csname figfile@ycb1-fig-suryanamaskar\endcsname{../workshop/images/new/ycb1_suryanamaskar_new.png}
\expandafter\def\csname figprompt@ycb1-fig-tadasan\endcsname{Posture: Tadasana. One adult practitioner (CS1) in the final held position of Tadasana, side view, standing, facing right. The whole body in frame with a small margin; head, spine, arms, hands, legs and feet positions exactly as in the notes block for this posture; weight and support points clear. No props unless the notes block names one. No labels. Aspect ratio 4:5 portrait.}
\expandafter\def\csname figfile@ycb1-fig-tadasan\endcsname{../workshop/images/new/ycb1_tadasan_new.png}
\expandafter\def\csname figprompt@ycb1-fig-teachers\endcsname{Do not generate portraits of real people. Replace with a plain timeline diagram of the teachers named in the checked notes block, or with photographs the author has permission to use.}
\expandafter\def\csname figfile@ycb1-fig-teachers\endcsname{../workshop/images/new/}
\expandafter\def\csname figprompt@ycb1-fig-teaching-demo\endcsname{Teaching scene: Essentials of Good Lesson Plan. A calm indoor yoga hall; an instructor in front (CS1 with a different top colour) and a small group of learners on mats, seen from behind or from the side, no recognisable faces, simple shapes. Shows only the idea in the frame title. Aspect ratio 16:9.}
\expandafter\def\csname figfile@ycb1-fig-teaching-demo\endcsname{../workshop/images/new/ycb1_teaching_demo_new.png}
\expandafter\def\csname figprompt@ycb1-fig-teaching-group\endcsname{Teaching scene: Principles of Teaching Yoga Protocol to Different Groups. A calm indoor yoga hall; an instructor in front (CS1 with a different top colour) and a small group of learners on mats, seen from behind or from the side, no recognisable faces, simple shapes. Shows only the idea in the frame title. Aspect ratio 16:9.}
\expandafter\def\csname figfile@ycb1-fig-teaching-group\endcsname{../workshop/images/new/ycb1_teaching_group_new.png}
\expandafter\def\csname figprompt@ycb1-fig-teaching-lecture\endcsname{Teaching scene: Conducting Yoga Practical Lessons: Precautions and Contraindications. A calm indoor yoga hall; an instructor in front (CS1 with a different top colour) and a small group of learners on mats, seen from behind or from the side, no recognisable faces, simple shapes. Shows only the idea in the frame title. Aspect ratio 16:9.}
\expandafter\def\csname figfile@ycb1-fig-teaching-lecture\endcsname{../workshop/images/new/ycb1_teaching_lecture_new.png}
\expandafter\def\csname figprompt@ycb1-fig-teaching-lesson\endcsname{Teaching scene: Models of Ideal Yoga Lesson Plans. A calm indoor yoga hall; an instructor in front (CS1 with a different top colour) and a small group of learners on mats, seen from behind or from the side, no recognisable faces, simple shapes. Shows only the idea in the frame title. Aspect ratio 16:9.}
\expandafter\def\csname figfile@ycb1-fig-teaching-lesson\endcsname{../workshop/images/new/ycb1_teaching_lesson_new.png}
\expandafter\def\csname figprompt@ycb1-fig-teaching-prep\endcsname{Teaching scene: Preparation for a Yoga Class (Before and During). A calm indoor yoga hall; an instructor in front (CS1 with a different top colour) and a small group of learners on mats, seen from behind or from the side, no recognisable faces, simple shapes. Shows only the idea in the frame title. Aspect ratio 16:9.}
\expandafter\def\csname figfile@ycb1-fig-teaching-prep\endcsname{../workshop/images/new/ycb1_teaching_prep_new.png}
\expandafter\def\csname figprompt@ycb1-fig-teaching-teacher\endcsname{A simple flat icon for Salient Features of an Ideal Yoga Instructor: one figure (CS1) with one symbol that stands for the idea, centred, plenty of white space. Aspect ratio 1:1.}
\expandafter\def\csname figfile@ycb1-fig-teaching-teacher\endcsname{../workshop/images/new/ycb1_teaching_teacher_new.pdf}
\expandafter\def\csname figprompt@ycb1-fig-threepartbreathing\endcsname{Preparatory breathing: Combination Breathing (Three-Part Breathing). One adult practitioner (CS1), side view, seated or lying as the notes block says, hands placed as the notes block says. Show the torso with a soft translucent highlight on the part that rises on the in-breath and a thin arrow for the direction of movement; a faint outline of the lungs and diaphragm may sit behind the highlight. Calm face, eyes closed. Aspect ratio 4:3 landscape.}
\expandafter\def\csname figfile@ycb1-fig-threepartbreathing\endcsname{../workshop/images/new/ycb1_threepartbreathing_new.png}
\expandafter\def\csname figprompt@ycb1-fig-tratak\endcsname{Practice: Trataka. One adult practitioner (CS1), seated, front three-quarter view, in the posture of the notes block. Show the key action in a small inset: the abdomen drawn in sharply for kapalabhati, the throat narrowed softly for ujjayi, the steady gaze on a candle flame at eye level for trataka, the tongue position for the mukh mudra. Use only the inset that matches Trataka. Aspect ratio 4:5 portrait.}
\expandafter\def\csname figfile@ycb1-fig-tratak\endcsname{../workshop/images/new/ycb1_tratak_new.png}
\expandafter\def\csname figprompt@ycb1-fig-trikonasan\endcsname{Posture: Trikonasana. One adult practitioner (CS1) in the final held position of Trikonasana, side view, standing, facing right. The whole body in frame with a small margin; head, spine, arms, hands, legs and feet positions exactly as in the notes block for this posture; weight and support points clear. No props unless the notes block names one. No labels. Aspect ratio 4:5 portrait.}
\expandafter\def\csname figfile@ycb1-fig-trikonasan\endcsname{../workshop/images/new/ycb1_trikonasan_new.png}
\expandafter\def\csname figprompt@ycb1-fig-ujjayi\endcsname{Practice: Ujjayi (Victorious Breath). One adult practitioner (CS1), seated, front three-quarter view, in the posture of the notes block. Show the key action in a small inset: the abdomen drawn in sharply for kapalabhati, the throat narrowed softly for ujjayi, the steady gaze on a candle flame at eye level for trataka, the tongue position for the mukh mudra. Use only the inset that matches Ujjayi (Victorious Breath). Aspect ratio 4:5 portrait.}
\expandafter\def\csname figfile@ycb1-fig-ujjayi\endcsname{../workshop/images/new/ycb1_ujjayi_new.png}
\expandafter\def\csname figprompt@ycb1-fig-ushtrasan\endcsname{Posture: Ushtrasana. One adult practitioner (CS1) in the final held position of Ushtrasana, front three-quarter view, seated on a plain mat. The whole body in frame with a small margin; head, spine, arms, hands, legs and feet positions exactly as in the notes block for this posture; weight and support points clear. No props unless the notes block names one. No labels. Aspect ratio 4:5 portrait.}
\expandafter\def\csname figfile@ycb1-fig-ushtrasan\endcsname{../workshop/images/new/ycb1_ushtrasan_new.png}
\expandafter\def\csname figprompt@ycb1-fig-uttanamandukasan\endcsname{Posture: Uttana Mandukasana. One adult practitioner (CS1) in the final held position of Uttana Mandukasana, front three-quarter view, seated on a plain mat. The whole body in frame with a small margin; head, spine, arms, hands, legs and feet positions exactly as in the notes block for this posture; weight and support points clear. No props unless the notes block names one. No labels. Aspect ratio 4:5 portrait.}
\expandafter\def\csname figfile@ycb1-fig-uttanamandukasan\endcsname{../workshop/images/new/ycb1_uttanamandukasan_new.png}
\expandafter\def\csname figprompt@ycb1-fig-uttanpadasan\endcsname{Posture: Uttanapadasana. One adult practitioner (CS1) in the final held position of Uttanapadasana, side view, lying on the back on a plain mat. The whole body in frame with a small margin; head, spine, arms, hands, legs and feet positions exactly as in the notes block for this posture; weight and support points clear. No props unless the notes block names one. No labels. Aspect ratio 3:2 landscape.}
\expandafter\def\csname figfile@ycb1-fig-uttanpadasan\endcsname{../workshop/images/new/ycb1_uttanpadasan_new.png}
\expandafter\def\csname figprompt@ycb1-fig-vajrasan\endcsname{Posture: Vajrasana. One adult practitioner (CS1) in the final held position of Vajrasana, front three-quarter view, seated on a plain mat. The whole body in frame with a small margin; head, spine, arms, hands, legs and feet positions exactly as in the notes block for this posture; weight and support points clear. No props unless the notes block names one. No labels. Aspect ratio 4:5 portrait.}
\expandafter\def\csname figfile@ycb1-fig-vajrasan\endcsname{../workshop/images/new/ycb1_vajrasan_new.png}
\expandafter\def\csname figprompt@ycb1-fig-vakrasan\endcsname{Posture: Vakrasana. One adult practitioner (CS1) in the final held position of Vakrasana, front three-quarter view, seated on a plain mat. The whole body in frame with a small margin; head, spine, arms, hands, legs and feet positions exactly as in the notes block for this posture; weight and support points clear. No props unless the notes block names one. No labels. Aspect ratio 4:5 portrait.}
\expandafter\def\csname figfile@ycb1-fig-vakrasan\endcsname{../workshop/images/new/ycb1_vakrasan_new.png}
\expandafter\def\csname figprompt@ycb1-fig-vipartikarni\endcsname{Posture: Vipareetakarani. One adult practitioner (CS1) in the final held position of Vipareetakarani, side view, lying on the back on a plain mat. The whole body in frame with a small margin; head, spine, arms, hands, legs and feet positions exactly as in the notes block for this posture; weight and support points clear. No props unless the notes block names one. No labels. Aspect ratio 3:2 landscape.}
\expandafter\def\csname figfile@ycb1-fig-vipartikarni\endcsname{../workshop/images/new/ycb1_vipartikarni_new.png}
\expandafter\def\csname figprompt@ycb1-fig-vrikshasan\endcsname{Posture: Vrikshasana. One adult practitioner (CS1) in the final held position of Vrikshasana, side view, standing, facing right. The whole body in frame with a small margin; head, spine, arms, hands, legs and feet positions exactly as in the notes block for this posture; weight and support points clear. No props unless the notes block names one. No labels. Aspect ratio 4:5 portrait.}
\expandafter\def\csname figfile@ycb1-fig-vrikshasan\endcsname{../workshop/images/new/ycb1_vrikshasan_new.png}
